% !TeX root = ../概率论与数理统计笔记.tex
% 由主文档通过 \input 载入；请编译主文档。
\setcounter{chapter}{0} % 保留原书章号。
\chapter{必要的数学工具}
\label{app:ps-prerequisites}

本附录按正文需求补充工具，每个主题以可直接使用的定义、条件、公式和短例题组织，不承担完整先修课程的替代作用。

\section{集合、计数与常用求和}
整理集合运算、可数集合、排列组合、二项式定理、几何级数与常见求和恒等式。
对无限级数区分绝对收敛与条件收敛，说明重排及逐项运算需要条件。

\section{微积分与积分交换}
回顾广义积分、多元积分、变量替换、泰勒展开及余项控制。
对非负或绝对可积函数说明交换求和与积分的适用情形；单调收敛、控制收敛与富比尼定理的概率应用列为进阶工具，并准确陈述使用条件。

\section{线性代数与二次型}
回顾矩阵秩、逆、迹、行列式、正交投影、正定与半正定二次型。
解释协方差矩阵与最小二乘问题使用这些工具的位置，补充实对称矩阵的正交对角化。

\section{特殊积分与渐近记号}
整理伽马函数、贝塔函数及其与常见密度归一化的关系，说明参数范围。
引入确定性及概率意义的阶记号时分别定义，配合极限定理说明近似项与余项的区别。






